% -*- coding: utf-8 -*-
% !TEX program = xelatex
\documentclass[plain,noanswer,medmath]{../../template/jnuexam}
\usepackage{../../template/kysx}

\begin{document}


\examtitle{1998}{1}


\makepart{填空题}{本题共5小题，每小题3分，满分15分}


\begin{problem}
$\lim_{x \to 0} \frac{\sqrt{1 + x} + \sqrt{1 - x} - 2}{x^2} =$ \fillin{}．
\end{problem}


\begin{problem}
设$z = \frac{1}{x}f(xy) + y\varphi(x + y)$，$f,\varphi$具有二阶连续导数，
则$\frac{\pd^2 z}{\pdx\pdy}=$ \fillin{}．
\end{problem}


\begin{problem}
设$L$为椭圆$\frac{x^2}{4} + \frac{y^2}{3} = 1$，其周长记为$a$，
则$\oint_L (2xy + 3x^2 + 4y^2)\ds =$ \fillin{}．
\end{problem}


\begin{problem}
设$A$为$n$阶矩阵，$|A| \ne 0$，$A^*$为$A$的伴随矩阵，$E$为$n$阶单位矩阵．
若$A$有特征值$\lambda$，则$(A^*)^2 + E$必有特征值 \fillin{}．
\end{problem}


\begin{problem}
设平面区域$D$由曲线$y = \frac{1}{x}$及直线$y = 0,x = 1,x = \e^2$所围成，
二维随机变量$(X,Y)$在区域$D$上服从均匀分布，则$(X,Y)$关于$X$的边缘概率密度在$x=2$处的值为 \fillin{}．
\end{problem}


\makepart{选择题}{本题共5小题，每小题3分，满分15分}


\begin{problem}
设$f(x)$连续，则$\frac{\d}{\dx}\int_0^x tf(x^2 - t^2) \dt =$\pickout{}
\begin{abcd}
\item $xf(x^2)$．
\item $- xf(x^2)$．
\item $2xf(x^2)$．
\item $- 2xf(x^2)$．
\end{abcd}
\end{problem}


\begin{problem}
函数$f(x) = (x^2 - x - 2)\left| x^3 - x \right|$不可导点的个数是\pickout{}
\begin{abcd}
\item $3$．
\item $2$．
\item $1$．
\item $0$．
\end{abcd}
\end{problem}


\begin{problem}
已知函数$y = y(x)$在任意点$x$处的增量$\Delta y = \frac{y\Delta x}{1 + x^2} + \alpha$，
且当$\Delta x \to 0$时，$\alpha$是$\Delta x$的高阶无穷小，$y(0) = \pi$，则$y(1)$等于\pickout{}
\begin{abcd}
\item $2\pi$．
\item $\pi$．
\item $e^{\frac{\pi}{4}}$．
\item $\pi \e^{\frac{\pi}{4}}$．
\end{abcd}
\end{problem}


\begin{problem}
设矩阵$\begin{pmatrix}
  a_1 & b_1 & c_1 \\
  a_2 & b_2 & c_2 \\
  a_3 & b_3 & c_3
\end{pmatrix}$是满秩的，
则直线$\frac{x - a_3}{a_1 - a_2} = \frac{y - b_3}{b_1 - b_2} = \frac{z - c_3}{c_1 - c_2}$
与直线$\frac{x - a_1}{a_2 - a_3} = \frac{y - b_1}{b_2 - b_3} = \frac{z - c_1}{c_2 - c_3}$\pickout{}
\begin{abcd*}
\item 相交于一点．
\item 重合．
\item 平行但不重合．
\item 异面．
\end{abcd*}
\end{problem}


\begin{problem}
设$A,B$是两个随机事件，且$0 < P(A) < 1,P(B) > 0,P(B|A) = P(B|\overline{A})$，则必有\pickout{}
\begin{abcd}
\item $P(A|B) = P(\overline{A} |B)$．
\item $P(A|B) \ne P(\overline{A} |B)$．
\item $P(AB) = P(A)P(B)$．
\item $P(AB) \ne P(A)P(B)$．
\end{abcd}
\end{problem}


\makepart{}{本题满分5分}

\begin{problem}
求直线$L:\frac{x - 1}{1} = \frac{y}{1} = \frac{z - 1}{-1}$
在平面$\Pi :x - y + 2z - 1 = 0$上的投影直线$L_0$的方程，
并求$L_0$绕$y$轴旋转一周所成曲面的方程．
\end{problem}


\makepart{}{本题满分6分}

\begin{problem}
确定常数$\lambda$，使在右半平面$x > 0$上的向量
$$A(x,y) = 2xy (x^4 + y^2)^{\lambda} i - x^2(x^4 + y^2)^{\lambda} j$$
为某二元函数$u(x,y)$的梯度，并求$u(x,y)$．
\end{problem}


\makepart{}{本题满分6分}

\begin{problem}
从船上向海中沉放某种探测仪器，按探测要求，需确定仪器的下沉深度$y$(从海平面算起)与下沉速度$v$
之间的函数关系．设仪器在重力作用下，从海平面由静止开始铅直下沉，在下沉过程中还受到阻力和浮力的作用．
设仪器的质量为$m$，体积为$B$，海水比重为$\rho$，仪器所受的阻力与下沉速度成正比，
比例系数为$k$（$k > 0$）．试建立$y$与$v$所满足的微分方程，并求出函数关系式$y = y(v)$．
\end{problem}


\makepart{}{本题满分7分}

\begin{problem}
计算$\iint_{\Sigma}\frac{ax\dy\dz + (z + a)^2 \dx\dy}{(x^2 + y^2 + z^2)^{\frac{1}{2}}}$，
其中$\Sigma$为下半球面$z = - \sqrt{a^2 - x^2 - y^2}$的上侧，$a$为大于零的常数．
\end{problem}


\makepart{}{本题满分6分}

\begin{problem}
求$\lim_{n \to \infty} \left(\frac{\sin \frac{\pi}{n}}{n + 1\vphantom{\frac{1}{1}}}
+ \frac{\sin \frac{2\pi}{n}}{n + \frac{1}{2}} + \cdots
+ \frac{\sin \pi\vphantom{\frac{1}{n}}}{n + \frac{1}{n}} \right)$．
\end{problem}


\makepart{}{本题满分5分}

\begin{problem}
设正项数列$\left\{a_n \right\}$单调减少，且$\sum_{n = 1}^{\infty}(- 1)^n a_n$发散，
试问级数$\sum_{n = 1}^{\infty}\left(\frac{1}{a_n + 1} \right)^n$是否收敛？并说明理由．
\end{problem}


\makepart{}{本题满分6分}

\begin{problem}
设$y = f(x)$是区间$[0,1]$上的任一非负连续函数．
\begin{enumerate}
\item 试证存在$x_0 \in(0,1)$，使得在区间$\left[0,x_0 \right]$上以$f(x_0)$为高的矩形面积，
      等于在区间$\left[x_0,1 \right]$上以$y = f(x)$为曲边的梯形面积．
\item 又设$f(x)$在区间$(0,1)$内可导，且$f'(x) > - \frac{2f(x)}{x}$，证明(I)中的$x_0$是唯一的．
\end{enumerate}
\end{problem}


\makepart{}{本题满分6分}

\begin{problem}
已知二次曲面方程$x^2 + ay^2 + z^2 + 2bxy + 2xz + 2yz = 4$，可以经过正交变换
$$\begin{pmatrix}
  x \\
  y \\
  z
\end{pmatrix} = P\begin{pmatrix}
  \xi \\
  \eta \\
  \zeta
\end{pmatrix}$$
化为椭圆柱面方程$\eta^2 + 4\zeta^2 = 4$，求$a,b$的值和正交矩阵$P$．
\end{problem}


\makepart{}{本题满分4分}

\begin{problem}
设$A$是$n$阶矩阵，若存在正整数$k$，使线性方程组$A^k x = 0$有解向量$\alpha$，
且$A^{k - 1}\alpha \ne 0$，证明：向量组$\alpha ,A\alpha , \cdots,A^{k - 1}\alpha$是线性无关的．
\end{problem}


\makepart{}{本题满分5分}

\begin{problem}
已知线性方程组
$$\text{①}\begin{cases}
  a_{11}x_1 + a_{12}x_2 + \cdots + a_{1,2n}x_{2n} = 0, \\
  a_{21}x_1 + a_{22}x_2 + \cdots + a_{2,2n}x_{2n} = 0, \\
  \cdots \cdots \cdots \cdots \cdots \cdots \cdots
  \cdots \cdots \cdots \cdots \cdots \cdots \\
  a_{n1}x_1 + a_{n2}x_2 + \cdots + a_{n,2n}x_{2n} = 0
\end{cases}$$
的一个基础解系为$(b_{11},b_{12},\cdots,b_{1,2n})^T, (b_{21},b_{22},\cdots,b_{2,2n})^T,
\cdots,(b_{n1},b_{n2},\cdots,b_{n,2n})^T$，试写出线性方程组
$$\text{②}\begin{cases}
  b_{11}y_1 + b_{12}y_2 + \cdots + b_{1,2n}y_{2n} = 0, \\
  b_{21}y_1 + b_{22}y_2 + \cdots + b_{2,2n}y_{2n} = 0, \\
  \cdots \cdots \cdots \cdots \cdots \cdots \cdots
  \cdots \cdots \cdots \cdots \cdots \cdots \\
  b_{n1}y_1 + b_{n2}y_2 + \cdots + b_{n,2n}y_{2n} = 0
\end{cases}$$
的通解，并说明理由．
\end{problem}


\makepart{}{本题满分6分}

\begin{problem}
设两个随机变量$X,Y$相互独立，且都服从均值为$0$、方差为$\frac{1}{2}$的正态分布，求随机变量$|X - Y|$的方差．
\end{problem}


\makepart{}{本题满分4分}

\begin{problem}
从正态总体$N\left(3.4,6^2 \right)$中抽取容量为$n$的样本，
如果要求其样本均值位于区间$(1.4,5.4)$内的概率不小于$0.95$，问样本容量$n$至少应取多大？
\par\centerline{附表：标准正态分布表
$\Phi(z) = \int_{-\infty}^z \frac{1}{\sqrt{2\pi}}\e^{-\frac{t^2}{2}}\dt$}\noindent
$$\begin{array}{|c|cccc|}
\hline
        z &  1.28 & 1.645 &  1.96 & 2.33 \\
\hline
  \Phi(z) & 0.900 & 0.950 & 0.975 & 0.990 \\
\hline
\end{array}$$
\end{problem}


\makepart{}{本题满分4分}

\begin{problem}
设某次考试的学生成绩服从正态分布，从中随机地抽取$36$位考生的成绩，
算得平均成绩为$66.5$分，标准差为$15$分．问在显著性水平$0.05$下，
是否可以认为这次考试全体考生的平均成绩为$70$分？并给出检验过程．
\par\centerline{附表：$t$分布表$P\{t(n) \le t_p(n)\} = p$}\noindent
$$\begin{array}{|c|cc|}
\hline
  \diagboxthree{n}{t_p(n)}{p} &   0.95 & 0.975 \\
\hline
                           35 & 1.6896 & 2.0301 \\
                           36 & 1.6883 & 2.0281 \\
\hline
\end{array}$$
\end{problem}


\end{document}
